\begin{exercise}{30}
Let $X$ be a topological space. Describe the category $h(\mathcal S(X))$.
Show that the endomorphisms of each object form a group. Which group is it?

\begin{proof}
The answer is the fundamental groupoid $\Pi_1(X)$: its objects
are points of $X$, its morphisms are homotopy classes relative
to endpoints of paths, and a composite traverses its right-hand
factor first.
By Definition~1.1.34, vertices and edges of $\mathcal S(X)$ are
precisely points and paths. Send an edge to its path class.
This respects Definition~1.2.5: constant paths give identities,
and in any singular triangle the broken path along edges
$(0,1),(1,2)$ is homotopic relative to endpoints to edge $(0,2)$.
Indeed, linearly interpolate these two paths inside the convex
topological triangle and then apply the singular simplex.
We obtain a functor $F\colon h(\mathcal S(X))\to\Pi_1(X)$
which is the identity on objects.

Every path class has a representative edge, so $F$ is full.
For faithfulness, first note that a pair of composable paths
$p,q$ is equal in $h(\mathcal S(X))$ to its concatenation
$p\cdot q$, where $p$ is traversed first. The singular triangle
\[
 (u_0,u_1,u_2)\longmapsto
 (p\cdot q)\bigl(\tfrac12u_1+u_2\bigr)
\]
has successive edges $p,q$ and long edge $p\cdot q$.
Consequently every morphism has a representative consisting
of one path, including identities represented by constant paths.

Suppose $H\colon[0,1]^2\to X$ is a homotopy relative to
endpoints from $p\colon x\to y$ to $p'\colon x\to y$.
The map
\[
 (t,s)\longmapsto(1-t,t(1-s),ts)
 \colon[0,1]^2\longrightarrow\Delta^2_{\mathrm{top}}
\]
is the quotient collapsing $\{0\}\times[0,1]$ to vertex $0$.
Since $H$ is constant there, it factors continuously through
a singular triangle. Its successive edges are $p$ and the
constant path at $y$, and its long edge is $p'$.
The triangle relation gives $[p']=[p]$ in $h(\mathcal S(X))$.
Thus $F$ is faithful and hence an isomorphism.

Reversing a path supplies its inverse: a path followed by its
reverse contracts relative to endpoints by shortening the
distance traversed. Hence all morphisms are invertible, and
the endomorphism group at $x$ is canonically the fundamental
group $\pi_1(X,x)$, with this concatenation convention.
\end{proof}
\end{exercise}
